% ----------------------------------------------------------------------
%  1.2.2. Пространственная решётка; 1.2.3. Кристаллографические символы;
%  1.2.4. Обратная решётка
% ----------------------------------------------------------------------
\subsection{Пространственная решётка}\label{ssec:lattice}

Любой узел решётки задаётся вектором
\begin{keyeq}\label{eq:Rm}
\vb R_m=m_1\vb a_1+m_2\vb a_2+m_3\vb a_3 ,
\end{keyeq}
где $\vb a_1$, $\vb a_2$, $\vb a_3$~— \term{векторы трансляций}, $m_1$, $m_2$,
$m_3$~— целые числа. Начало координат можно поместить в любой узел.
Решётку можно получить бесконечным повторением параллелепипеда,
построенного на кратчайших некомпланарных трансляциях,~— \term{элементарной
ячейки}. Если на её гранях и внутри нет узлов, ячейку называют
\term{примитивной}: каждый её вершинный узел принадлежит восьми соседним
ячейкам, и на ячейку приходится ровно один узел.

\begin{figure}[ht]
\centering
\begin{tikzpicture}[x={(1cm,0cm)}, y={(0.55cm,0.38cm)}, z={(0.12cm,1cm)}, scale=1.35]
  \coordinate (O) at (0,0,0);
  \coordinate (A) at (2.2,0,0); \coordinate (B) at (0.35,1.7,0); \coordinate (C) at (0.25,0.2,1.8);
  \draw[atomLine!70, dashed] (O) -- (B) -- ($(A)+(B)$) ($(B)$) -- ($(B)+(C)$);
  \draw[atomLine] (A) -- ($(A)+(B)$) -- ($(A)+(B)+(C)$) -- ($(A)+(C)$) -- cycle
    ($(C)$) -- ($(B)+(C)$) -- ($(A)+(B)+(C)$) (C) -- ($(A)+(C)$) (O) -- (C);
  \draw[bvec, fmBlue] (O) -- (A) node[below, font=\small, black]{$\vb a_1=\vb a$};
  \draw[bvec, fmBlue] (O) -- (B) node[above left=-2pt, font=\small, black]{$\vb a_2=\vb b$};
  \draw[bvec, fmBlue] (O) -- (C) node[left, font=\small, black]{$\vb a_3=\vb c$};
  \node[font=\small, fmRed] at (0.42,0.32,0.08) {$\gamma$};
  \node[font=\small, fmRed] at (0.55,0.0,0.45) {$\beta$};
  \node[font=\small, fmRed] at (0.12,0.45,0.55) {$\alpha$};
\end{tikzpicture}
\caption{Элементарная ячейка и её параметры: длины рёбер $a$, $b$, $c$ и
углы $\alpha=\widehat{\vb b\vb c}$, $\beta=\widehat{\vb a\vb c}$,
$\gamma=\widehat{\vb a\vb b}$}
\label{fig:cell}
\end{figure}

Длины $a$, $b$, $c$ векторов трансляций и углы $\alpha$, $\beta$, $\gamma$ между
ними~— \term{параметры решётки} (рис.~\ref{fig:cell}). Кристаллографическая
система координат строится на этих векторах; в общем случае она
\emph{косоугольная} с разными масштабами по осям. По симметрии
расположения узлов решётки делятся на семь \term{сингоний}
(табл.~\ref{tab:syngonies}).

\begin{table}[ht]
\centering
\caption{Сингонии (кристаллографические системы координат)}\label{tab:syngonies}
\small
\begin{tabular}{@{}lcl@{}}
\toprule
Сингония & Число параметров & Параметры \\
\midrule
Триклинная & 6 & $a\ne b\ne c$;\quad $\alpha\ne\beta\ne\gamma\ne90^\circ$ \\
Моноклинная & 4 & $a\ne b\ne c$;\quad $\alpha=\gamma=90^\circ$, $\beta\ne90^\circ$ \\
Ромбическая & 3 & $a\ne b\ne c$;\quad $\alpha=\beta=\gamma=90^\circ$ \\
Тетрагональная & 2 & $a=b\ne c$;\quad $\alpha=\beta=\gamma=90^\circ$ \\
Ромбоэдрическая (тригональная) & 2 & $a=b=c$;\quad $\alpha=\beta=\gamma\ne90^\circ$ \\
Гексагональная & 2 & $a=b\ne c$;\quad $\alpha=\beta=90^\circ$, $\gamma=120^\circ$ \\
Кубическая & 1 & $a=b=c$;\quad $\alpha=\beta=\gamma=90^\circ$ \\
\bottomrule
\end{tabular}
\end{table}

\paragraph{Метрическая матрица}
Все метрические свойства решётки (длины, углы, объёмы) содержатся в
\term{метрической матрице} $G$, составленной из скалярных произведений
базисных векторов:
\begin{equation}\label{eq:G}
g_{ij}=(\vb a_i\vb a_j),\qquad a_i=\sqrt{g_{ii}},\qquad
\cos\alpha_{ij}=\frac{g_{ij}}{\sqrt{g_{ii}g_{jj}}},\qquad V=\sqrt{\det G} .
\end{equation}
Матрица симметрична ($g_{ij}=g_{ji}$). Для важнейших сингоний
\[
G_\text{куб}=\begin{pmatrix}a^2&0&0\\0&a^2&0\\0&0&a^2\end{pmatrix},\quad
G_\text{ромб}=\begin{pmatrix}a^2&0&0\\0&b^2&0\\0&0&c^2\end{pmatrix},\quad
G_\text{гекс}=\begin{pmatrix}a^2&-a^2/2&0\\-a^2/2&a^2&0\\0&0&c^2\end{pmatrix}
\]
(тетрагональная~— как ромбическая с $b=a$). Пользу матрицы $G$ мы увидим в
п.~\ref{ssec:formulas}: с её помощью любые углы и расстояния в
косоугольной системе вычисляются по одной формуле.

\paragraph{Сложные ячейки}
Ради наглядности (чтобы ячейка отражала симметрию решётки) часто
выбирают не примитивные, а \emph{сложные} ячейки, у которых есть
дополнительные узлы в центре объёма или граней (правила выбора~—
п.~1.3.6). Положение узла тогда задаётся суммой $\vb R_m+\vb R_n$, где
$\vb R_n=n_1\vb a_1+n_2\vb a_2+n_3\vb a_3$ с \emph{дробными} $n_i$~— это
\term{базис} решётки:
\begin{itemize}
\item объёмноцентрированная ($I$): $[\![000]\!]$, $[\![\tfrac12\tfrac12\tfrac12]\!]$;
\item гранецентрированная ($F$): $[\![000]\!]$, $[\![\tfrac12\tfrac120]\!]$,
$[\![0\tfrac12\tfrac12]\!]$, $[\![\tfrac120\tfrac12]\!]$;
\item базоцентрированная ($A$, $B$ или $C$~— по центрируемой грани), например
$A$: $[\![000]\!]$, $[\![0\tfrac12\tfrac12]\!]$.
\end{itemize}

\subsection{Кристаллографические символы}\label{ssec:symbols}

\paragraph{Символ узла}
Узел с вектором $\vb R=m\vb a+n\vb b+p\vb c$ обозначают $[\![mnp]\!]$. Цифры
пишут подряд без запятых (запятые~— только для двузначных индексов) и
читают порознь; минус ставят \emph{над} цифрой: $[\![1\br20]\!]$ читается
«один, минус два, ноль». У узлов сложной ячейки индексы дробные: центр
объёма~— $[\![\tfrac12\tfrac12\tfrac12]\!]$, центр грани~—
$[\![0\tfrac12\tfrac12]\!]$.

\paragraph{Символ направления}
Решётка~— семейство параллельных узловых прямых. Его характеризует
прямая, проходящая через начало координат; её символ $[uvw]$~— индексы
ближайшего к началу узла на этой прямой (целые взаимно простые числа).
Оси координат имеют символы $[100]$, $[010]$, $[001]$ \emph{в любой
сингонии}~— одно из главных удобств кристаллографической символики.
Совокупность направлений, связанных симметрией, обозначают угловыми
скобками: в кубической сингонии $\nfam{100}$~— это $[100]$, $[010]$,
$[001]$, $[\br100]$, $[0\br10]$, $[00\br1]$, т.\,е. направления,
отличающиеся перестановкой и знаками индексов.

\begin{figure}[ht]
\centering
\begin{subfigure}[b]{0.42\textwidth}\centering
\begin{tikzpicture}[cubeproj, scale=2.4]
  \CubeEdges
  \foreach \v/\l/\a in {{(1,0,0)}/{[100]}/north, {(0,1,0)}/{[010]}/south east, {(0,0,1)}/{[001]}/east,
      {(1,1,0)}/{[110]}/west, {(1,0,1)}/{[101]}/south west, {(1,1,1)}/{[111]}/south}{
    \draw[bvec, fmBlue, line width=0.9pt] (0,0,0) -- \v;
    \node[font=\scriptsize, anchor=\a, fill=white, inner sep=1pt, fill opacity=0.8, text opacity=1] at \v {$\l$};}
  \draw[bvec, fmRed, line width=0.9pt] (1,0,0) -- (0,1,0);
  \node[font=\scriptsize, fmRed, anchor=north] at (0.5,0.45,0) {$[\br110]$};
  \fill (0,0,0) circle (0.6pt);
\end{tikzpicture}
\caption{}
\end{subfigure}\hfill
\begin{subfigure}[b]{0.55\textwidth}\centering
\tikzset{zone100/.style={fmBlue, line width=0.6pt},
  zone110/.style={fmBlue!50, line width=0.4pt},
  pole100/.style={rectangle, fill=fmRed, inner sep=0pt, minimum size=4pt},
  pole110/.style={ellipse, fill=fmGreen, inner sep=0pt, minimum width=4.5pt, minimum height=3pt},
  pole111/.style={regular polygon, regular polygon sides=3, fill=fmAmber!90!black, inner sep=0pt, minimum size=6pt},
  poleH/.style={circle, fill=atomLine, inner sep=0pt, minimum size=2pt}}
\begin{tikzpicture}[scale=3.0]
  \draw[atomLine, line width=0.8pt] (0,0) circle (1);
  % создано tools/gen_stereo.py: стандартная проекция (001) куба (r = 1)
\draw[zone100] (0.0000,-1.0000) -- (0.0000,-0.9742) -- (0.0000,-0.9490) -- (0.0000,-0.9244) -- (0.0000,-0.9004) -- (0.0000,-0.8770) -- (0.0000,-0.8541) -- (0.0000,-0.8317) -- (0.0000,-0.8098) -- (0.0000,-0.7883) -- (0.0000,-0.7673) -- (0.0000,-0.7467) -- (0.0000,-0.7265) -- (0.0000,-0.7067) -- (0.0000,-0.6873) -- (0.0000,-0.6682) -- (0.0000,-0.6494) -- (0.0000,-0.6310) -- (0.0000,-0.6128) -- (0.0000,-0.5949) -- (0.0000,-0.5774) -- (0.0000,-0.5600) -- (0.0000,-0.5430) -- (0.0000,-0.5261) -- (0.0000,-0.5095) -- (0.0000,-0.4931) -- (0.0000,-0.4770) -- (0.0000,-0.4610) -- (0.0000,-0.4452) -- (0.0000,-0.4296) -- (0.0000,-0.4142) -- (0.0000,-0.3990) -- (0.0000,-0.3839) -- (0.0000,-0.3689) -- (0.0000,-0.3541) -- (0.0000,-0.3395) -- (0.0000,-0.3249) -- (0.0000,-0.3105) -- (0.0000,-0.2962) -- (0.0000,-0.2820) -- (0.0000,-0.2679) -- (0.0000,-0.2540) -- (0.0000,-0.2401) -- (0.0000,-0.2263) -- (0.0000,-0.2126) -- (0.0000,-0.1989) -- (0.0000,-0.1853) -- (0.0000,-0.1718) -- (0.0000,-0.1584) -- (0.0000,-0.1450) -- (0.0000,-0.1317) -- (0.0000,-0.1184) -- (0.0000,-0.1051) -- (0.0000,-0.0919) -- (0.0000,-0.0787) -- (0.0000,-0.0655) -- (0.0000,-0.0524) -- (0.0000,-0.0393) -- (0.0000,-0.0262) -- (0.0000,-0.0131) -- (0.0000,-0.0000) -- (-0.0000,0.0131) -- (-0.0000,0.0262) -- (-0.0000,0.0393) -- (-0.0000,0.0524) -- (-0.0000,0.0655) -- (-0.0000,0.0787) -- (-0.0000,0.0919) -- (-0.0000,0.1051) -- (-0.0000,0.1184) -- (-0.0000,0.1317) -- (-0.0000,0.1450) -- (-0.0000,0.1584) -- (-0.0000,0.1718) -- (-0.0000,0.1853) -- (-0.0000,0.1989) -- (-0.0000,0.2126) -- (-0.0000,0.2263) -- (-0.0000,0.2401) -- (-0.0000,0.2540) -- (-0.0000,0.2679) -- (-0.0000,0.2820) -- (-0.0000,0.2962) -- (-0.0000,0.3105) -- (-0.0000,0.3249) -- (-0.0000,0.3395) -- (-0.0000,0.3541) -- (-0.0000,0.3689) -- (-0.0000,0.3839) -- (-0.0000,0.3990) -- (-0.0000,0.4142) -- (-0.0000,0.4296) -- (-0.0000,0.4452) -- (-0.0000,0.4610) -- (-0.0000,0.4770) -- (-0.0000,0.4931) -- (-0.0000,0.5095) -- (-0.0000,0.5261) -- (-0.0000,0.5430) -- (-0.0000,0.5600) -- (-0.0000,0.5774) -- (-0.0000,0.5949) -- (-0.0000,0.6128) -- (-0.0000,0.6310) -- (-0.0000,0.6494) -- (-0.0000,0.6682) -- (-0.0000,0.6873) -- (-0.0000,0.7067) -- (-0.0000,0.7265) -- (-0.0000,0.7467) -- (-0.0000,0.7673) -- (-0.0000,0.7883) -- (-0.0000,0.8098) -- (-0.0000,0.8317) -- (-0.0000,0.8541) -- (-0.0000,0.8770) -- (-0.0000,0.9004) -- (-0.0000,0.9244) -- (-0.0000,0.9490) -- (-0.0000,0.9742) -- (-0.0000,1.0000);
\draw[zone100] (1.0000,0.0000) -- (0.9742,0.0000) -- (0.9490,0.0000) -- (0.9244,0.0000) -- (0.9004,0.0000) -- (0.8770,0.0000) -- (0.8541,0.0000) -- (0.8317,0.0000) -- (0.8098,0.0000) -- (0.7883,0.0000) -- (0.7673,0.0000) -- (0.7467,0.0000) -- (0.7265,0.0000) -- (0.7067,0.0000) -- (0.6873,0.0000) -- (0.6682,0.0000) -- (0.6494,0.0000) -- (0.6310,0.0000) -- (0.6128,0.0000) -- (0.5949,0.0000) -- (0.5774,0.0000) -- (0.5600,0.0000) -- (0.5430,0.0000) -- (0.5261,0.0000) -- (0.5095,0.0000) -- (0.4931,0.0000) -- (0.4770,0.0000) -- (0.4610,0.0000) -- (0.4452,0.0000) -- (0.4296,0.0000) -- (0.4142,0.0000) -- (0.3990,0.0000) -- (0.3839,0.0000) -- (0.3689,0.0000) -- (0.3541,0.0000) -- (0.3395,0.0000) -- (0.3249,0.0000) -- (0.3105,0.0000) -- (0.2962,0.0000) -- (0.2820,0.0000) -- (0.2679,0.0000) -- (0.2540,0.0000) -- (0.2401,0.0000) -- (0.2263,0.0000) -- (0.2126,0.0000) -- (0.1989,0.0000) -- (0.1853,0.0000) -- (0.1718,0.0000) -- (0.1584,0.0000) -- (0.1450,0.0000) -- (0.1317,0.0000) -- (0.1184,0.0000) -- (0.1051,0.0000) -- (0.0919,0.0000) -- (0.0787,0.0000) -- (0.0655,0.0000) -- (0.0524,0.0000) -- (0.0393,0.0000) -- (0.0262,0.0000) -- (0.0131,0.0000) -- (0.0000,0.0000) -- (-0.0131,-0.0000) -- (-0.0262,-0.0000) -- (-0.0393,-0.0000) -- (-0.0524,-0.0000) -- (-0.0655,-0.0000) -- (-0.0787,-0.0000) -- (-0.0919,-0.0000) -- (-0.1051,-0.0000) -- (-0.1184,-0.0000) -- (-0.1317,-0.0000) -- (-0.1450,-0.0000) -- (-0.1584,-0.0000) -- (-0.1718,-0.0000) -- (-0.1853,-0.0000) -- (-0.1989,-0.0000) -- (-0.2126,-0.0000) -- (-0.2263,-0.0000) -- (-0.2401,-0.0000) -- (-0.2540,-0.0000) -- (-0.2679,-0.0000) -- (-0.2820,-0.0000) -- (-0.2962,-0.0000) -- (-0.3105,-0.0000) -- (-0.3249,-0.0000) -- (-0.3395,-0.0000) -- (-0.3541,-0.0000) -- (-0.3689,-0.0000) -- (-0.3839,-0.0000) -- (-0.3990,-0.0000) -- (-0.4142,-0.0000) -- (-0.4296,-0.0000) -- (-0.4452,-0.0000) -- (-0.4610,-0.0000) -- (-0.4770,-0.0000) -- (-0.4931,-0.0000) -- (-0.5095,-0.0000) -- (-0.5261,-0.0000) -- (-0.5430,-0.0000) -- (-0.5600,-0.0000) -- (-0.5774,-0.0000) -- (-0.5949,-0.0000) -- (-0.6128,-0.0000) -- (-0.6310,-0.0000) -- (-0.6494,-0.0000) -- (-0.6682,-0.0000) -- (-0.6873,-0.0000) -- (-0.7067,-0.0000) -- (-0.7265,-0.0000) -- (-0.7467,-0.0000) -- (-0.7673,-0.0000) -- (-0.7883,-0.0000) -- (-0.8098,-0.0000) -- (-0.8317,-0.0000) -- (-0.8541,-0.0000) -- (-0.8770,-0.0000) -- (-0.9004,-0.0000) -- (-0.9244,-0.0000) -- (-0.9490,-0.0000) -- (-0.9742,-0.0000) -- (-1.0000,-0.0000);
\draw[zone110] (0.7071,-0.7071) -- (0.6888,-0.6888) -- (0.6710,-0.6710) -- (0.6536,-0.6536) -- (0.6367,-0.6367) -- (0.6201,-0.6201) -- (0.6039,-0.6039) -- (0.5881,-0.5881) -- (0.5726,-0.5726) -- (0.5574,-0.5574) -- (0.5426,-0.5426) -- (0.5280,-0.5280) -- (0.5137,-0.5137) -- (0.4997,-0.4997) -- (0.4860,-0.4860) -- (0.4725,-0.4725) -- (0.4592,-0.4592) -- (0.4462,-0.4462) -- (0.4333,-0.4333) -- (0.4207,-0.4207) -- (0.4082,-0.4082) -- (0.3960,-0.3960) -- (0.3839,-0.3839) -- (0.3720,-0.3720) -- (0.3603,-0.3603) -- (0.3487,-0.3487) -- (0.3373,-0.3373) -- (0.3260,-0.3260) -- (0.3148,-0.3148) -- (0.3038,-0.3038) -- (0.2929,-0.2929) -- (0.2821,-0.2821) -- (0.2714,-0.2714) -- (0.2609,-0.2609) -- (0.2504,-0.2504) -- (0.2400,-0.2400) -- (0.2298,-0.2298) -- (0.2196,-0.2196) -- (0.2095,-0.2095) -- (0.1994,-0.1994) -- (0.1895,-0.1895) -- (0.1796,-0.1796) -- (0.1698,-0.1698) -- (0.1600,-0.1600) -- (0.1503,-0.1503) -- (0.1407,-0.1407) -- (0.1311,-0.1311) -- (0.1215,-0.1215) -- (0.1120,-0.1120) -- (0.1025,-0.1025) -- (0.0931,-0.0931) -- (0.0837,-0.0837) -- (0.0743,-0.0743) -- (0.0650,-0.0650) -- (0.0557,-0.0557) -- (0.0463,-0.0463) -- (0.0371,-0.0371) -- (0.0278,-0.0278) -- (0.0185,-0.0185) -- (0.0093,-0.0093) -- (0.0000,-0.0000) -- (-0.0093,0.0093) -- (-0.0185,0.0185) -- (-0.0278,0.0278) -- (-0.0371,0.0371) -- (-0.0463,0.0463) -- (-0.0557,0.0557) -- (-0.0650,0.0650) -- (-0.0743,0.0743) -- (-0.0837,0.0837) -- (-0.0931,0.0931) -- (-0.1025,0.1025) -- (-0.1120,0.1120) -- (-0.1215,0.1215) -- (-0.1311,0.1311) -- (-0.1407,0.1407) -- (-0.1503,0.1503) -- (-0.1600,0.1600) -- (-0.1698,0.1698) -- (-0.1796,0.1796) -- (-0.1895,0.1895) -- (-0.1994,0.1994) -- (-0.2095,0.2095) -- (-0.2196,0.2196) -- (-0.2298,0.2298) -- (-0.2400,0.2400) -- (-0.2504,0.2504) -- (-0.2609,0.2609) -- (-0.2714,0.2714) -- (-0.2821,0.2821) -- (-0.2929,0.2929) -- (-0.3038,0.3038) -- (-0.3148,0.3148) -- (-0.3260,0.3260) -- (-0.3373,0.3373) -- (-0.3487,0.3487) -- (-0.3603,0.3603) -- (-0.3720,0.3720) -- (-0.3839,0.3839) -- (-0.3960,0.3960) -- (-0.4082,0.4082) -- (-0.4207,0.4207) -- (-0.4333,0.4333) -- (-0.4462,0.4462) -- (-0.4592,0.4592) -- (-0.4725,0.4725) -- (-0.4860,0.4860) -- (-0.4997,0.4997) -- (-0.5137,0.5137) -- (-0.5280,0.5280) -- (-0.5426,0.5426) -- (-0.5574,0.5574) -- (-0.5726,0.5726) -- (-0.5881,0.5881) -- (-0.6039,0.6039) -- (-0.6201,0.6201) -- (-0.6367,0.6367) -- (-0.6536,0.6536) -- (-0.6710,0.6710) -- (-0.6888,0.6888) -- (-0.7071,0.7071);
\draw[zone110] (-0.7071,-0.7071) -- (-0.6888,-0.6888) -- (-0.6710,-0.6710) -- (-0.6536,-0.6536) -- (-0.6367,-0.6367) -- (-0.6201,-0.6201) -- (-0.6039,-0.6039) -- (-0.5881,-0.5881) -- (-0.5726,-0.5726) -- (-0.5574,-0.5574) -- (-0.5426,-0.5426) -- (-0.5280,-0.5280) -- (-0.5137,-0.5137) -- (-0.4997,-0.4997) -- (-0.4860,-0.4860) -- (-0.4725,-0.4725) -- (-0.4592,-0.4592) -- (-0.4462,-0.4462) -- (-0.4333,-0.4333) -- (-0.4207,-0.4207) -- (-0.4082,-0.4082) -- (-0.3960,-0.3960) -- (-0.3839,-0.3839) -- (-0.3720,-0.3720) -- (-0.3603,-0.3603) -- (-0.3487,-0.3487) -- (-0.3373,-0.3373) -- (-0.3260,-0.3260) -- (-0.3148,-0.3148) -- (-0.3038,-0.3038) -- (-0.2929,-0.2929) -- (-0.2821,-0.2821) -- (-0.2714,-0.2714) -- (-0.2609,-0.2609) -- (-0.2504,-0.2504) -- (-0.2400,-0.2400) -- (-0.2298,-0.2298) -- (-0.2196,-0.2196) -- (-0.2095,-0.2095) -- (-0.1994,-0.1994) -- (-0.1895,-0.1895) -- (-0.1796,-0.1796) -- (-0.1698,-0.1698) -- (-0.1600,-0.1600) -- (-0.1503,-0.1503) -- (-0.1407,-0.1407) -- (-0.1311,-0.1311) -- (-0.1215,-0.1215) -- (-0.1120,-0.1120) -- (-0.1025,-0.1025) -- (-0.0931,-0.0931) -- (-0.0837,-0.0837) -- (-0.0743,-0.0743) -- (-0.0650,-0.0650) -- (-0.0557,-0.0557) -- (-0.0463,-0.0463) -- (-0.0371,-0.0371) -- (-0.0278,-0.0278) -- (-0.0185,-0.0185) -- (-0.0093,-0.0093) -- (-0.0000,-0.0000) -- (0.0093,0.0093) -- (0.0185,0.0185) -- (0.0278,0.0278) -- (0.0371,0.0371) -- (0.0463,0.0463) -- (0.0557,0.0557) -- (0.0650,0.0650) -- (0.0743,0.0743) -- (0.0837,0.0837) -- (0.0931,0.0931) -- (0.1025,0.1025) -- (0.1120,0.1120) -- (0.1215,0.1215) -- (0.1311,0.1311) -- (0.1407,0.1407) -- (0.1503,0.1503) -- (0.1600,0.1600) -- (0.1698,0.1698) -- (0.1796,0.1796) -- (0.1895,0.1895) -- (0.1994,0.1994) -- (0.2095,0.2095) -- (0.2196,0.2196) -- (0.2298,0.2298) -- (0.2400,0.2400) -- (0.2504,0.2504) -- (0.2609,0.2609) -- (0.2714,0.2714) -- (0.2821,0.2821) -- (0.2929,0.2929) -- (0.3038,0.3038) -- (0.3148,0.3148) -- (0.3260,0.3260) -- (0.3373,0.3373) -- (0.3487,0.3487) -- (0.3603,0.3603) -- (0.3720,0.3720) -- (0.3839,0.3839) -- (0.3960,0.3960) -- (0.4082,0.4082) -- (0.4207,0.4207) -- (0.4333,0.4333) -- (0.4462,0.4462) -- (0.4592,0.4592) -- (0.4725,0.4725) -- (0.4860,0.4860) -- (0.4997,0.4997) -- (0.5137,0.5137) -- (0.5280,0.5280) -- (0.5426,0.5426) -- (0.5574,0.5574) -- (0.5726,0.5726) -- (0.5881,0.5881) -- (0.6039,0.6039) -- (0.6201,0.6201) -- (0.6367,0.6367) -- (0.6536,0.6536) -- (0.6710,0.6710) -- (0.6888,0.6888) -- (0.7071,0.7071);
\draw[zone110] (0.0000,-1.0000) -- (-0.0182,-0.9815) -- (-0.0357,-0.9630) -- (-0.0526,-0.9445) -- (-0.0688,-0.9261) -- (-0.0845,-0.9077) -- (-0.0996,-0.8893) -- (-0.1142,-0.8710) -- (-0.1282,-0.8528) -- (-0.1417,-0.8346) -- (-0.1547,-0.8165) -- (-0.1672,-0.7985) -- (-0.1793,-0.7805) -- (-0.1910,-0.7626) -- (-0.2022,-0.7448) -- (-0.2130,-0.7271) -- (-0.2234,-0.7095) -- (-0.2334,-0.6919) -- (-0.2430,-0.6745) -- (-0.2523,-0.6571) -- (-0.2612,-0.6398) -- (-0.2698,-0.6226) -- (-0.2780,-0.6055) -- (-0.2860,-0.5884) -- (-0.2936,-0.5715) -- (-0.3009,-0.5546) -- (-0.3080,-0.5378) -- (-0.3147,-0.5211) -- (-0.3212,-0.5045) -- (-0.3274,-0.4879) -- (-0.3333,-0.4714) -- (-0.3390,-0.4550) -- (-0.3445,-0.4386) -- (-0.3497,-0.4224) -- (-0.3546,-0.4061) -- (-0.3594,-0.3900) -- (-0.3639,-0.3739) -- (-0.3682,-0.3579) -- (-0.3723,-0.3419) -- (-0.3761,-0.3260) -- (-0.3798,-0.3101) -- (-0.3833,-0.2943) -- (-0.3865,-0.2785) -- (-0.3896,-0.2628) -- (-0.3925,-0.2471) -- (-0.3951,-0.2315) -- (-0.3976,-0.2159) -- (-0.4000,-0.2003) -- (-0.4021,-0.1848) -- (-0.4040,-0.1693) -- (-0.4058,-0.1538) -- (-0.4074,-0.1383) -- (-0.4089,-0.1229) -- (-0.4101,-0.1075) -- (-0.4112,-0.0921) -- (-0.4121,-0.0767) -- (-0.4129,-0.0614) -- (-0.4135,-0.0460) -- (-0.4139,-0.0307) -- (-0.4141,-0.0153) -- (-0.4142,-0.0000) -- (-0.4141,0.0153) -- (-0.4139,0.0307) -- (-0.4135,0.0460) -- (-0.4129,0.0614) -- (-0.4121,0.0767) -- (-0.4112,0.0921) -- (-0.4101,0.1075) -- (-0.4089,0.1229) -- (-0.4074,0.1383) -- (-0.4058,0.1538) -- (-0.4040,0.1693) -- (-0.4021,0.1848) -- (-0.4000,0.2003) -- (-0.3976,0.2159) -- (-0.3951,0.2315) -- (-0.3925,0.2471) -- (-0.3896,0.2628) -- (-0.3865,0.2785) -- (-0.3833,0.2943) -- (-0.3798,0.3101) -- (-0.3761,0.3260) -- (-0.3723,0.3419) -- (-0.3682,0.3579) -- (-0.3639,0.3739) -- (-0.3594,0.3900) -- (-0.3546,0.4061) -- (-0.3497,0.4224) -- (-0.3445,0.4386) -- (-0.3390,0.4550) -- (-0.3333,0.4714) -- (-0.3274,0.4879) -- (-0.3212,0.5045) -- (-0.3147,0.5211) -- (-0.3080,0.5378) -- (-0.3009,0.5546) -- (-0.2936,0.5715) -- (-0.2860,0.5884) -- (-0.2780,0.6055) -- (-0.2698,0.6226) -- (-0.2612,0.6398) -- (-0.2523,0.6571) -- (-0.2430,0.6745) -- (-0.2334,0.6919) -- (-0.2234,0.7095) -- (-0.2130,0.7271) -- (-0.2022,0.7448) -- (-0.1910,0.7626) -- (-0.1793,0.7805) -- (-0.1672,0.7985) -- (-0.1547,0.8165) -- (-0.1417,0.8346) -- (-0.1282,0.8528) -- (-0.1142,0.8710) -- (-0.0996,0.8893) -- (-0.0845,0.9077) -- (-0.0688,0.9261) -- (-0.0526,0.9445) -- (-0.0357,0.9630) -- (-0.0182,0.9815) -- (-0.0000,1.0000);
\draw[zone110] (0.0000,-1.0000) -- (0.0182,-0.9815) -- (0.0357,-0.9630) -- (0.0526,-0.9445) -- (0.0688,-0.9261) -- (0.0845,-0.9077) -- (0.0996,-0.8893) -- (0.1142,-0.8710) -- (0.1282,-0.8528) -- (0.1417,-0.8346) -- (0.1547,-0.8165) -- (0.1672,-0.7985) -- (0.1793,-0.7805) -- (0.1910,-0.7626) -- (0.2022,-0.7448) -- (0.2130,-0.7271) -- (0.2234,-0.7095) -- (0.2334,-0.6919) -- (0.2430,-0.6745) -- (0.2523,-0.6571) -- (0.2612,-0.6398) -- (0.2698,-0.6226) -- (0.2780,-0.6055) -- (0.2860,-0.5884) -- (0.2936,-0.5715) -- (0.3009,-0.5546) -- (0.3080,-0.5378) -- (0.3147,-0.5211) -- (0.3212,-0.5045) -- (0.3274,-0.4879) -- (0.3333,-0.4714) -- (0.3390,-0.4550) -- (0.3445,-0.4386) -- (0.3497,-0.4224) -- (0.3546,-0.4061) -- (0.3594,-0.3900) -- (0.3639,-0.3739) -- (0.3682,-0.3579) -- (0.3723,-0.3419) -- (0.3761,-0.3260) -- (0.3798,-0.3101) -- (0.3833,-0.2943) -- (0.3865,-0.2785) -- (0.3896,-0.2628) -- (0.3925,-0.2471) -- (0.3951,-0.2315) -- (0.3976,-0.2159) -- (0.4000,-0.2003) -- (0.4021,-0.1848) -- (0.4040,-0.1693) -- (0.4058,-0.1538) -- (0.4074,-0.1383) -- (0.4089,-0.1229) -- (0.4101,-0.1075) -- (0.4112,-0.0921) -- (0.4121,-0.0767) -- (0.4129,-0.0614) -- (0.4135,-0.0460) -- (0.4139,-0.0307) -- (0.4141,-0.0153) -- (0.4142,-0.0000) -- (0.4141,0.0153) -- (0.4139,0.0307) -- (0.4135,0.0460) -- (0.4129,0.0614) -- (0.4121,0.0767) -- (0.4112,0.0921) -- (0.4101,0.1075) -- (0.4089,0.1229) -- (0.4074,0.1383) -- (0.4058,0.1538) -- (0.4040,0.1693) -- (0.4021,0.1848) -- (0.4000,0.2003) -- (0.3976,0.2159) -- (0.3951,0.2315) -- (0.3925,0.2471) -- (0.3896,0.2628) -- (0.3865,0.2785) -- (0.3833,0.2943) -- (0.3798,0.3101) -- (0.3761,0.3260) -- (0.3723,0.3419) -- (0.3682,0.3579) -- (0.3639,0.3739) -- (0.3594,0.3900) -- (0.3546,0.4061) -- (0.3497,0.4224) -- (0.3445,0.4386) -- (0.3390,0.4550) -- (0.3333,0.4714) -- (0.3274,0.4879) -- (0.3212,0.5045) -- (0.3147,0.5211) -- (0.3080,0.5378) -- (0.3009,0.5546) -- (0.2936,0.5715) -- (0.2860,0.5884) -- (0.2780,0.6055) -- (0.2698,0.6226) -- (0.2612,0.6398) -- (0.2523,0.6571) -- (0.2430,0.6745) -- (0.2334,0.6919) -- (0.2234,0.7095) -- (0.2130,0.7271) -- (0.2022,0.7448) -- (0.1910,0.7626) -- (0.1793,0.7805) -- (0.1672,0.7985) -- (0.1547,0.8165) -- (0.1417,0.8346) -- (0.1282,0.8528) -- (0.1142,0.8710) -- (0.0996,0.8893) -- (0.0845,0.9077) -- (0.0688,0.9261) -- (0.0526,0.9445) -- (0.0357,0.9630) -- (0.0182,0.9815) -- (0.0000,1.0000);
\draw[zone110] (1.0000,0.0000) -- (0.9815,-0.0182) -- (0.9630,-0.0357) -- (0.9445,-0.0526) -- (0.9261,-0.0688) -- (0.9077,-0.0845) -- (0.8893,-0.0996) -- (0.8710,-0.1142) -- (0.8528,-0.1282) -- (0.8346,-0.1417) -- (0.8165,-0.1547) -- (0.7985,-0.1672) -- (0.7805,-0.1793) -- (0.7626,-0.1910) -- (0.7448,-0.2022) -- (0.7271,-0.2130) -- (0.7095,-0.2234) -- (0.6919,-0.2334) -- (0.6745,-0.2430) -- (0.6571,-0.2523) -- (0.6398,-0.2612) -- (0.6226,-0.2698) -- (0.6055,-0.2780) -- (0.5884,-0.2860) -- (0.5715,-0.2936) -- (0.5546,-0.3009) -- (0.5378,-0.3080) -- (0.5211,-0.3147) -- (0.5045,-0.3212) -- (0.4879,-0.3274) -- (0.4714,-0.3333) -- (0.4550,-0.3390) -- (0.4386,-0.3445) -- (0.4224,-0.3497) -- (0.4061,-0.3546) -- (0.3900,-0.3594) -- (0.3739,-0.3639) -- (0.3579,-0.3682) -- (0.3419,-0.3723) -- (0.3260,-0.3761) -- (0.3101,-0.3798) -- (0.2943,-0.3833) -- (0.2785,-0.3865) -- (0.2628,-0.3896) -- (0.2471,-0.3925) -- (0.2315,-0.3951) -- (0.2159,-0.3976) -- (0.2003,-0.4000) -- (0.1848,-0.4021) -- (0.1693,-0.4040) -- (0.1538,-0.4058) -- (0.1383,-0.4074) -- (0.1229,-0.4089) -- (0.1075,-0.4101) -- (0.0921,-0.4112) -- (0.0767,-0.4121) -- (0.0614,-0.4129) -- (0.0460,-0.4135) -- (0.0307,-0.4139) -- (0.0153,-0.4141) -- (0.0000,-0.4142) -- (-0.0153,-0.4141) -- (-0.0307,-0.4139) -- (-0.0460,-0.4135) -- (-0.0614,-0.4129) -- (-0.0767,-0.4121) -- (-0.0921,-0.4112) -- (-0.1075,-0.4101) -- (-0.1229,-0.4089) -- (-0.1383,-0.4074) -- (-0.1538,-0.4058) -- (-0.1693,-0.4040) -- (-0.1848,-0.4021) -- (-0.2003,-0.4000) -- (-0.2159,-0.3976) -- (-0.2315,-0.3951) -- (-0.2471,-0.3925) -- (-0.2628,-0.3896) -- (-0.2785,-0.3865) -- (-0.2943,-0.3833) -- (-0.3101,-0.3798) -- (-0.3260,-0.3761) -- (-0.3419,-0.3723) -- (-0.3579,-0.3682) -- (-0.3739,-0.3639) -- (-0.3900,-0.3594) -- (-0.4061,-0.3546) -- (-0.4224,-0.3497) -- (-0.4386,-0.3445) -- (-0.4550,-0.3390) -- (-0.4714,-0.3333) -- (-0.4879,-0.3274) -- (-0.5045,-0.3212) -- (-0.5211,-0.3147) -- (-0.5378,-0.3080) -- (-0.5546,-0.3009) -- (-0.5715,-0.2936) -- (-0.5884,-0.2860) -- (-0.6055,-0.2780) -- (-0.6226,-0.2698) -- (-0.6398,-0.2612) -- (-0.6571,-0.2523) -- (-0.6745,-0.2430) -- (-0.6919,-0.2334) -- (-0.7095,-0.2234) -- (-0.7271,-0.2130) -- (-0.7448,-0.2022) -- (-0.7626,-0.1910) -- (-0.7805,-0.1793) -- (-0.7985,-0.1672) -- (-0.8165,-0.1547) -- (-0.8346,-0.1417) -- (-0.8528,-0.1282) -- (-0.8710,-0.1142) -- (-0.8893,-0.0996) -- (-0.9077,-0.0845) -- (-0.9261,-0.0688) -- (-0.9445,-0.0526) -- (-0.9630,-0.0357) -- (-0.9815,-0.0182) -- (-1.0000,-0.0000);
\draw[zone110] (1.0000,0.0000) -- (0.9815,0.0182) -- (0.9630,0.0357) -- (0.9445,0.0526) -- (0.9261,0.0688) -- (0.9077,0.0845) -- (0.8893,0.0996) -- (0.8710,0.1142) -- (0.8528,0.1282) -- (0.8346,0.1417) -- (0.8165,0.1547) -- (0.7985,0.1672) -- (0.7805,0.1793) -- (0.7626,0.1910) -- (0.7448,0.2022) -- (0.7271,0.2130) -- (0.7095,0.2234) -- (0.6919,0.2334) -- (0.6745,0.2430) -- (0.6571,0.2523) -- (0.6398,0.2612) -- (0.6226,0.2698) -- (0.6055,0.2780) -- (0.5884,0.2860) -- (0.5715,0.2936) -- (0.5546,0.3009) -- (0.5378,0.3080) -- (0.5211,0.3147) -- (0.5045,0.3212) -- (0.4879,0.3274) -- (0.4714,0.3333) -- (0.4550,0.3390) -- (0.4386,0.3445) -- (0.4224,0.3497) -- (0.4061,0.3546) -- (0.3900,0.3594) -- (0.3739,0.3639) -- (0.3579,0.3682) -- (0.3419,0.3723) -- (0.3260,0.3761) -- (0.3101,0.3798) -- (0.2943,0.3833) -- (0.2785,0.3865) -- (0.2628,0.3896) -- (0.2471,0.3925) -- (0.2315,0.3951) -- (0.2159,0.3976) -- (0.2003,0.4000) -- (0.1848,0.4021) -- (0.1693,0.4040) -- (0.1538,0.4058) -- (0.1383,0.4074) -- (0.1229,0.4089) -- (0.1075,0.4101) -- (0.0921,0.4112) -- (0.0767,0.4121) -- (0.0614,0.4129) -- (0.0460,0.4135) -- (0.0307,0.4139) -- (0.0153,0.4141) -- (0.0000,0.4142) -- (-0.0153,0.4141) -- (-0.0307,0.4139) -- (-0.0460,0.4135) -- (-0.0614,0.4129) -- (-0.0767,0.4121) -- (-0.0921,0.4112) -- (-0.1075,0.4101) -- (-0.1229,0.4089) -- (-0.1383,0.4074) -- (-0.1538,0.4058) -- (-0.1693,0.4040) -- (-0.1848,0.4021) -- (-0.2003,0.4000) -- (-0.2159,0.3976) -- (-0.2315,0.3951) -- (-0.2471,0.3925) -- (-0.2628,0.3896) -- (-0.2785,0.3865) -- (-0.2943,0.3833) -- (-0.3101,0.3798) -- (-0.3260,0.3761) -- (-0.3419,0.3723) -- (-0.3579,0.3682) -- (-0.3739,0.3639) -- (-0.3900,0.3594) -- (-0.4061,0.3546) -- (-0.4224,0.3497) -- (-0.4386,0.3445) -- (-0.4550,0.3390) -- (-0.4714,0.3333) -- (-0.4879,0.3274) -- (-0.5045,0.3212) -- (-0.5211,0.3147) -- (-0.5378,0.3080) -- (-0.5546,0.3009) -- (-0.5715,0.2936) -- (-0.5884,0.2860) -- (-0.6055,0.2780) -- (-0.6226,0.2698) -- (-0.6398,0.2612) -- (-0.6571,0.2523) -- (-0.6745,0.2430) -- (-0.6919,0.2334) -- (-0.7095,0.2234) -- (-0.7271,0.2130) -- (-0.7448,0.2022) -- (-0.7626,0.1910) -- (-0.7805,0.1793) -- (-0.7985,0.1672) -- (-0.8165,0.1547) -- (-0.8346,0.1417) -- (-0.8528,0.1282) -- (-0.8710,0.1142) -- (-0.8893,0.0996) -- (-0.9077,0.0845) -- (-0.9261,0.0688) -- (-0.9445,0.0526) -- (-0.9630,0.0357) -- (-0.9815,0.0182) -- (-1.0000,0.0000);
\node[pole100] at (1.0000,0.0000) {};
\node[font=\scriptsize, inner sep=1pt, anchor=180] at (1.0000,0.0000) {$100$};
\node[pole100] at (-1.0000,0.0000) {};
\node[font=\scriptsize, inner sep=1pt, anchor=360] at (-1.0000,0.0000) {$\br{1}00$};
\node[pole100] at (0.0000,0.0000) {};
\node[font=\scriptsize, inner sep=1pt, anchor=90] at (0.0000,0.0000) {$001$};
\node[pole100] at (0.0000,1.0000) {};
\node[font=\scriptsize, inner sep=1pt, anchor=270] at (0.0000,1.0000) {$010$};
\node[pole100] at (0.0000,-1.0000) {};
\node[font=\scriptsize, inner sep=1pt, anchor=90] at (0.0000,-1.0000) {$0\br{1}0$};
\node[pole110] at (0.4142,0.0000) {};
\node[font=\scriptsize, inner sep=1pt, anchor=180] at (0.4142,0.0000) {$101$};
\node[pole110] at (-0.4142,0.0000) {};
\node[font=\scriptsize, inner sep=1pt, anchor=360] at (-0.4142,0.0000) {$\br{1}01$};
\node[pole110] at (0.7071,0.7071) {};
\node[font=\scriptsize, inner sep=1pt, anchor=225] at (0.7071,0.7071) {$110$};
\node[pole110] at (0.7071,-0.7071) {};
\node[font=\scriptsize, inner sep=1pt, anchor=135] at (0.7071,-0.7071) {$1\br{1}0$};
\node[pole110] at (-0.7071,0.7071) {};
\node[font=\scriptsize, inner sep=1pt, anchor=315] at (-0.7071,0.7071) {$\br{1}10$};
\node[pole110] at (-0.7071,-0.7071) {};
\node[font=\scriptsize, inner sep=1pt, anchor=45] at (-0.7071,-0.7071) {$\br{1}\br{1}0$};
\node[pole110] at (0.0000,0.4142) {};
\node[font=\scriptsize, inner sep=1pt, anchor=270] at (0.0000,0.4142) {$011$};
\node[pole110] at (0.0000,-0.4142) {};
\node[font=\scriptsize, inner sep=1pt, anchor=90] at (0.0000,-0.4142) {$0\br{1}1$};
\node[pole111] at (0.3660,0.3660) {};
\node[font=\scriptsize, inner sep=1pt, anchor=225] at (0.3660,0.3660) {$111$};
\node[pole111] at (0.3660,-0.3660) {};
\node[font=\scriptsize, inner sep=1pt, anchor=135] at (0.3660,-0.3660) {$1\br{1}1$};
\node[pole111] at (-0.3660,0.3660) {};
\node[font=\scriptsize, inner sep=1pt, anchor=315] at (-0.3660,0.3660) {$\br{1}11$};
\node[pole111] at (-0.3660,-0.3660) {};
\node[font=\scriptsize, inner sep=1pt, anchor=45] at (-0.3660,-0.3660) {$\br{1}\br{1}1$};
\node[poleH] at (0.0000,0.6180) {};
\node[poleH] at (0.0000,-0.6180) {};
\node[poleH] at (0.4472,0.8944) {};
\node[poleH] at (0.4472,-0.8944) {};
\node[poleH] at (-0.4472,0.8944) {};
\node[poleH] at (-0.4472,-0.8944) {};
\node[poleH] at (0.8944,0.4472) {};
\node[poleH] at (0.8944,-0.4472) {};
\node[poleH] at (-0.8944,0.4472) {};
\node[poleH] at (-0.8944,-0.4472) {};
\node[poleH] at (0.6180,0.0000) {};
\node[poleH] at (-0.6180,0.0000) {};
\node[poleH] at (0.0000,0.2361) {};
\node[poleH] at (0.0000,-0.2361) {};
\node[poleH] at (0.2361,0.0000) {};
\node[poleH] at (-0.2361,0.0000) {};
\node[poleH] at (0.2899,0.5798) {};
\node[poleH] at (0.2899,-0.5798) {};
\node[poleH] at (-0.2899,0.5798) {};
\node[poleH] at (-0.2899,-0.5798) {};
\node[poleH] at (0.5798,0.2899) {};
\node[poleH] at (0.5798,-0.2899) {};
\node[poleH] at (-0.5798,0.2899) {};
\node[poleH] at (-0.5798,-0.2899) {};
\node[poleH] at (0.2247,0.2247) {};
\node[poleH] at (0.2247,-0.2247) {};
\node[poleH] at (-0.2247,0.2247) {};
\node[poleH] at (-0.2247,-0.2247) {};
\node[poleH] at (0.2000,0.4000) {};
\node[poleH] at (0.2000,-0.4000) {};
\node[poleH] at (-0.2000,0.4000) {};
\node[poleH] at (-0.2000,-0.4000) {};
\node[poleH] at (0.4000,0.2000) {};
\node[poleH] at (0.4000,-0.2000) {};
\node[poleH] at (-0.4000,0.2000) {};
\node[poleH] at (-0.4000,-0.2000) {};
\node[poleH] at (0.5000,0.5000) {};
\node[poleH] at (0.5000,-0.5000) {};
\node[poleH] at (-0.5000,0.5000) {};
\node[poleH] at (-0.5000,-0.5000) {};

\end{tikzpicture}
\caption{}
\end{subfigure}
\caption{Важнейшие направления кубического кристалла: \textit{а}~— в
элементарной ячейке (начало координат~— в левой нижней передней
вершине); \textit{б}~— стандартная стереографическая проекция $(001)$
(расчёт): полюсы $\{100\}$ (квадраты), $\{110\}$ (овалы), $\{111\}$
(треугольники), $\{210\}$, $\{211\}$, $\{221\}$ (точки) и следы
плоскостей $\{100\}$ и $\{110\}$. В кубическом кристалле направление
$[uvw]$ перпендикулярно плоскости $(uvw)$, поэтому полюсы плоскостей и
выходы направлений совпадают}
\label{fig:cubic-dirs}
\end{figure}

\paragraph{Примитивные ячейки ГЦК и ОЦК}
В сложных ячейках можно выделить примитивные, построенные на кратчайших
трансляциях (рис.~\ref{fig:prim}). Для ГЦК-решётки это векторы
$\frac a2[110]$, $\frac a2[011]$, $\frac a2[101]$, для ОЦК~—
$\frac a2[11\br1]$, $\frac a2[\br111]$, $\frac a2[1\br11]$. В матричной
записи (строки~— векторы)
\begin{equation}\label{eq:prim}
A_\text{ГЦК}=\frac a2\begin{pmatrix}1&1&0\\0&1&1\\1&0&1\end{pmatrix},\qquad
A_\text{ОЦК}=\frac a2\begin{pmatrix}1&1&\br1\\\br1&1&1\\1&\br1&1\end{pmatrix}.
\end{equation}
Объём примитивной ячейки ГЦК равен $a^3/4$, ОЦК~— $a^3/2$ (в
кубической ячейке 4 и 2 узла соответственно).

\begin{figure}[ht]
\centering
\begin{subfigure}[b]{0.45\textwidth}\centering
\begin{tikzpicture}[cubeproj, scale=2.4]
  \CubeEdges[atomLine!60]
  \coordinate (v1) at (0.5,0.5,0); \coordinate (v2) at (0,0.5,0.5); \coordinate (v3) at (0.5,0,0.5);
  \draw[fmRed!70, line width=0.5pt] (v1) -- (0.5,1,0.5) -- (v2) (v1) -- (1,0.5,0.5) -- (v3)
    (v2) -- (0.5,0.5,1) -- (v3) (0.5,1,0.5) -- (1,1,1) (1,0.5,0.5) -- (1,1,1) (0.5,0.5,1) -- (1,1,1);
  \FCCback[2.2mm] \FCCmid[2.2mm] \FCCfront[2.2mm]
  \draw[bvec, fmRed] (0,0,0) -- (v1); \draw[bvec, fmRed] (0,0,0) -- (v2); \draw[bvec, fmRed] (0,0,0) -- (v3);
\end{tikzpicture}
\caption{ГЦК}
\end{subfigure}\hfill
\begin{subfigure}[b]{0.45\textwidth}\centering
\begin{tikzpicture}[cubeproj, scale=2.4]
  \CubeEdges[atomLine!60]
  \foreach \p in {(0,0,0),(1,0,0),(0,1,0),(1,1,0),(0,0,1),(1,0,1),(0,1,1),(1,1,1)}{\node[atom=2.6mm] at \p {};}
  \node[atomB=2.6mm] at (0.5,0.5,0.5) {};
  \draw[bvec, fmRed] (0.5,0.5,0.5) -- (1,1,0);
  \draw[bvec, fmRed] (0.5,0.5,0.5) -- (0,1,1);
  \draw[bvec, fmRed] (0.5,0.5,0.5) -- (1,0,1);
\end{tikzpicture}
\caption{ОЦК}
\end{subfigure}
\caption{Векторы примитивных трансляций (красные стрелки) в ГЦК- и
ОЦК-решётках. Для ГЦК показана и сама примитивная ячейка~—
ромбоэдр с углом $60^\circ$; для ОЦК векторы проведены из центрального узла
(ромбоэдр с углом $109{,}5^\circ$ выходит за пределы куба)}
\label{fig:prim}
\end{figure}

\paragraph{Символ плоскости (индексы Миллера)}
Решётку можно представить и как семейство параллельных узловых
плоскостей. Пусть плоскость отсекает на осях отрезки $A$, $B$, $C$; её
уравнение $x/A+y/B+z/C=1$ выполняется для узлов $x=ma$, $y=nb$, $z=pc$, т.\,е.
$m(a/A)+n(b/B)+p(c/C)=1$. Отсюда отношения $a/A$, $b/B$, $c/C$ рациональны и
их можно записать через три взаимно простых целых числа:
\begin{keyeq}\label{eq:miller}
h:k:l=\frac aA:\frac bB:\frac cC .
\end{keyeq}

\begin{opredelenie}[Индексы Миллера]
Три взаимно простых целых числа $h$, $k$, $l$, обратно пропорциональные
отрезкам, которые плоскость отсекает на осях (в единицах $a$, $b$, $c$).
Символ плоскости записывают в круглых скобках: $(hkl)$.
\end{opredelenie}

\begin{shagi}
\item Найти отрезки, отсекаемые плоскостью на осях, в единицах периодов
(если плоскость параллельна оси, отрезок~— $\infty$; если проходит через
начало~— сдвинуть начало в соседний узел).
\item Взять обратные величины.
\item Привести к наименьшим целым числам; отрицательный индекс отметить
чертой сверху.
\end{shagi}

\begin{primer}[Индексы плоскости]
Плоскость отсекает отрезки $1$, $2$ и $\infty$. Обратные величины $1$,
$\frac12$, $0$; умножая на 2, получаем $(210)$. Плоскость с отрезками
$\frac12$, $1$, $\infty$ имеет те же индексы $(210)$ и параллельна первой:
все плоскости семейства имеют один символ.
\end{primer}

Подставив $h=q\,a/A$, $k=q\,b/B$, $l=q\,c/C$ в уравнение плоскости, получим
\begin{equation}\label{eq:plane-q}
hm+kn+lp=q ,
\end{equation}
где $q$~— целое число, «номер» плоскости в семействе: $q=0$~— плоскость
через начало координат, $q=1$~— ближайшая к нему; она отсекает отрезки
$a/h$, $b/k$, $c/l$. Расположение важнейших плоскостей кубической решётки
показано на рис.~\ref{fig:planes}.

\begin{figure}[ht]
\centering
\tikzset{plshade/.style={fill=fmBlue!25, draw=fmBlue, line width=0.6pt, fill opacity=0.8}}
\foreach \pl/\lab in {
  {(1,0,0) -- (1,1,0) -- (1,1,1) -- (1,0,1) -- cycle}/{(100)},
  {(1,0,0) -- (0,1,0) -- (0,1,1) -- (1,0,1) -- cycle}/{(110)},
  {(1,0,0) -- (0,1,0) -- (0,0,1) -- cycle}/{(111)},
  {(0.5,0,0) -- (0,1,0) -- (0,1,1) -- (0.5,0,1) -- cycle}/{(210)}}{%
\begin{tikzpicture}[cubeproj, scale=1.6, baseline=(b)]
  \coordinate (b) at (0,0,0);
  \CubeEdges
  \filldraw[plshade] \pl;
  \node[font=\small] at (0.6,0,-0.32) {$\lab$};
\end{tikzpicture}\hspace{4mm}}
\caption{Плоскости $(100)$, $(110)$, $(111)$ и $(210)$ в кубической ячейке
(начало координат~— в левой нижней передней вершине)}
\label{fig:planes}
\end{figure}

\paragraph{Кристаллическая форма и множитель повторяемости}
Совокупность плоскостей, связанных элементами симметрии, обозначают
фигурными скобками $\{hkl\}$ и называют \term{кристаллической формой}. В
кубической сингонии $\{100\}$~— это $(100)$, $(010)$, $(001)$, $(\br100)$,
$(0\br10)$, $(00\br1)$~— шесть граней куба. Число плоскостей формы~—
\term{множитель повторяемости} $p$: для кубической сингонии
$p_{\{100\}}=6$, $p_{\{110\}}=12$, $p_{\{111\}}=8$, $p_{\{123\}}=48$. Все плоскости
одной формы имеют одинаковое межплоскостное расстояние.

\paragraph{Четырёхиндексные символы в гексагональной системе}
Базисная плоскость гексагональной решётки имеет три равноправные оси
$\vb a_1$, $\vb a_2$, $\vb a_3=-(\vb a_1+\vb a_2)$ под углами $120^\circ$
(рис.~\ref{fig:hex4}). Поэтому плоскости обозначают четырьмя индексами
$(hkil)$, где третий индекс относится к оси $\vb a_3$. Он не независим:
\begin{keyeq}\label{eq:hkil}
i=-(h+k) .
\end{keyeq}
(Вывод: если след плоскости отсекает на осях $\vb a_1$, $\vb a_2$, $\vb a_3$
отрезки $A_1$, $A_2$, $-A_3$, то из подобия треугольников
$1/A_1+1/A_2=1/A_3$, а индексы обратны отрезкам.) Выгода~— в том, что
симметрично равнозначные плоскости получаются перестановкой первых трёх
индексов: шесть граней призмы первого рода~— $(10\br10)$, $(01\br10)$,
$(\br1100)$, $(\br1010)$, $(0\br110)$, $(1\br100)$, т.\,е. форма $\{10\br10\}$, тогда
как в трёх индексах они выглядят непохоже: $(100)$, $(010)$, $(\br110)$\ldots{}
Вычисления обычно ведут в трёх индексах, а третий индекс заменяют точкой:
$(hk.l)$.

Направления тоже записывают четырьмя индексами $[uvtw]$ с условием
$u+v+t=0$. Если в трёх индексах $\vb R=U\vb a_1+V\vb a_2+W\vb c$, то, подставив
$\vb a_3=-(\vb a_1+\vb a_2)$ в $\vb R=u\vb a_1+v\vb a_2+t\vb a_3+w\vb c$, находим
$U=2u+v$, $V=u+2v$, $W=w$, откуда (с точностью до общего множителя)
\begin{keyeq}\label{eq:uvtw}
u=2U-V,\qquad v=2V-U,\qquad t=-(u+v),\qquad w=3W .
\end{keyeq}

\begin{primer}[Переход к четырём индексам]
$[100]\to[2\br1\br10]$; $[110]\to[11\br20]$; $[001]\to[0001]$;
$[\br110]\to[\br1100]$. Направления $\nfam{11\br20}$~— это шесть направлений
вдоль осей $\pm\vb a_1$, $\pm\vb a_2$, $\pm\vb a_3$ (плотноупакованные ряды в
ГПУ-металлах).
\end{primer}

\begin{figure}[ht]
\centering
\begin{tikzpicture}[scale=1.25]
  \foreach \i in {-2,...,2}{\foreach \j in {-2,...,2}{
    \pgfmathsetmacro{\px}{\i+0.5*\j}\pgfmathsetmacro{\py}{0.866*\j}
    \pgfmathparse{abs(\i+\j)<=2 ? 1 : 0}\ifnum\pgfmathresult=1
      \fill[atomA, draw=atomLine, line width=0.3pt] (\px,\py) circle (0.07);\fi}}
  \draw[fmGray, thin] (0:1) \foreach \a in {60,120,...,300}{-- (\a:1)} -- cycle;
  \draw[bvec, fmBlue] (0,0) -- (1,0) node[below right=-2pt, font=\small, black]{$\vb a_1$};
  \draw[bvec, fmBlue] (0,0) -- (120:1) node[above left=-2pt, font=\small, black]{$\vb a_2$};
  \draw[bvec, fmBlue] (0,0) -- (240:1) node[below left=-2pt, font=\small, black]{$\vb a_3$};
  \draw[fmRed, very thick] (1,0) -- (60:1);
  \node[fmRed, font=\scriptsize, anchor=west] at (30:0.95) {$(10\br10)$};
  \draw[fmGreen!70!black, very thick] ($(1,0)!-0.25!(120:1)$) -- ($(1,0)!1.25!(120:1)$);
  \node[fmGreen!70!black, font=\scriptsize, anchor=south east] at ($(1,0)!1.25!(120:1)$) {$(11\br20)$};
  \draw[bvec, fmAmber!90!black] (0,0) -- (60:1.6) node[above right=-2pt, font=\scriptsize]{$[11\br20]$};
  \node[font=\scriptsize, fmBlue, anchor=north] at (1.65,-0.05) {$[2\br1\br10]$};
\end{tikzpicture}
\caption{Гексагональная система: оси $\vb a_1$, $\vb a_2$, $\vb a_3$ в
базисной плоскости, следы плоскостей призмы первого рода $(10\br10)$ и
второго рода $(11\br20)$, направления $[2\br1\br10]$ (вдоль $\vb a_1$) и
$[11\br20]$ (вдоль $\vb a_1+\vb a_2$)}
\label{fig:hex4}
\end{figure}

\subsection{Обратная решётка}\label{ssec:reciprocal}

\begin{opredelenie}[Обратная решётка]
Решётка, базисные векторы которой $\vb b_1$, $\vb b_2$, $\vb b_3$ (их обозначают
также $\vb a^*$, $\vb b^*$, $\vb c^*$) связаны с базисом прямой решётки
условием
\begin{equation}\label{eq:recip-def}
(\vb a_i\vb b_j)=\delta_{ij}\qquad(\delta_{ii}=1,\ \delta_{ij}=0\text{ при }i\ne j).
\end{equation}
\end{opredelenie}
Из определения $\vb b_1\perp\vb a_2,\vb a_3$, т.\,е. $\vb b_1\parallel[\vb a_2\times\vb a_3]$;
нормировка $(\vb a_1\vb b_1)=1$ даёт
\begin{keyeq}\label{eq:b1}
\vb b_1=\frac{[\vb a_2\times\vb a_3]}{V},\quad
\vb b_2=\frac{[\vb a_3\times\vb a_1]}{V},\quad
\vb b_3=\frac{[\vb a_1\times\vb a_2]}{V},\qquad V=(\vb a_1[\vb a_2\times\vb a_3]) .
\end{keyeq}
Решётки взаимны: прямая решётка~— обратная к обратной,
$\vb a_1=[\vb b_2\times\vb b_3]/V^*$, и объёмы ячеек связаны условием $VV^*=1$. В
физике твёрдого тела обычно используют определение
$(\vb a_i\vb b_j)=2\pi\delta_{ij}$ (так в гл.~3)~— все векторы обратной решётки
просто умножаются на $2\pi$.

Метрическая матрица обратной решётки $g^{ij}=(\vb b_i\vb b_j)$ равна $G^{-1}$:
\[
G^{-1}_\text{куб}=\frac1{a^2}I,\qquad
G^{-1}_\text{ромб}=\operatorname{diag}\Bigl(\frac1{a^2},\frac1{b^2},\frac1{c^2}\Bigr),\qquad
G^{-1}_\text{гекс}=\begin{pmatrix}\frac4{3a^2}&\frac2{3a^2}&0\\\frac2{3a^2}&\frac4{3a^2}&0\\0&0&\frac1{c^2}\end{pmatrix}.
\]

\begin{podrobnee}[Матричная запись и взаимность ГЦК--ОЦК]
Если векторы прямого базиса записать строками матрицы $A$, а векторы
обратного~— столбцами матрицы $B$, условие~\eqref{eq:recip-def} принимает
вид $AB=I$, т.\,е. $B=A^{-1}$. Для ГЦК-решётки из~\eqref{eq:prim}
\[
A_\text{ГЦК}^{-1}=\frac1a\begin{pmatrix}1&\br1&1\\1&1&\br1\\\br1&1&1\end{pmatrix},
\]
столбцы которой~— векторы $\frac1a(1,1,\br1)$, $\frac1a(\br1,1,1)$,
$\frac1a(1,\br1,1)$: это примитивные векторы \emph{ОЦК}-решётки с
кубической ячейкой $2/a$. Итак, обратная к ГЦК-решётке с ребром $a$~—
ОЦК-решётка с ребром $2/a$ (объём кубической ячейки $8/a^3$), и наоборот.
Метрические матрицы связаны: $\vb b=G^{-1}\vb a$, $\vb a=G\vb b$, $GG^{-1}=I$.
\end{podrobnee}

\paragraph{Вектор обратной решётки и его свойства}
Узел обратной решётки задаётся вектором
$\vb H_{HKL}=H\vb b_1+K\vb b_2+L\vb b_3$, где $H=qh$, $K=qk$, $L=ql$ (общий
множитель $q$ может быть больше единицы).

\begin{keyeq}\label{eq:H-props}
\vb H_{hkl}=h\vb b_1+k\vb b_2+l\vb b_3\ \perp\ (hkl),\qquad
|\vb H_{hkl}|=\frac1{d_{hkl}} .
\end{keyeq}
\begin{itemize}
\item \emph{Перпендикулярность.} Для любого узлового вектора
$\vb R=m\vb a_1+n\vb a_2+p\vb a_3$, лежащего в плоскости $(hkl)$, проходящей через
начало, $(\vb H\vb R)=hm+kn+lp=0$ по~\eqref{eq:plane-q} с $q=0$.
\item \emph{Модуль.} Ближайшая к началу плоскость семейства отсекает на
оси $x$ отрезок $\vb a_1/h$; его проекция на единичную нормаль
$\vb H/|\vb H|$~— межплоскостное расстояние:
$d_{hkl}=\frac{\vb a_1}{h}\cdot\frac{h\vb b_1+k\vb b_2+l\vb b_3}{|\vb H|}=\frac1{|\vb H|}$.
\end{itemize}
Так \vazhno{каждому семейству плоскостей прямой решётки отвечает узел
обратной}: направление на него~— нормаль к плоскостям, расстояние~—
обратное межплоскостное (рис.~\ref{fig:H-vector}).

\begin{figure}[ht]
\centering
\begin{tikzpicture}[scale=1.15]
  \begin{scope}
    \clip (-1.6,-1.0) rectangle (4.6,3.0);
    \foreach \q in {-2,...,5}{
      \draw[fmBlue!60] ($\q*(1.2,0)+(-3*0.8,3*1.0)$) -- ($\q*(1.2,0)+(4*0.8,-4*1.0)$);}
    \foreach \m in {-2,...,5}{\foreach \n in {-2,...,4}{
      \fill[atomA, draw=atomLine, line width=0.3pt] ($\m*(1.2,0)+\n*(0.4,1.0)$) circle (0.06);}}
  \end{scope}
  \draw[bvec, atomLine] (0,0) -- (1.2,0) node[below, font=\small]{$\vb a$};
  \draw[bvec, atomLine] (0,0) -- (0.4,1.0) node[left, font=\small]{$\vb b$};
  \draw[bvec, fmRed, line width=1.1pt] (0,0) -- ($1.6*(0.781,0.625)$) node[above, font=\small]{$\vb H_{11}$};
  \draw[dimline, fmGreen!60!black] ($(0,0)+0.0*(0.781,0.625)+(0.39,-0.49)$) -- ($0.937*(0.781,0.625)+(0.39,-0.49)$);
  \node[font=\small, fmGreen!60!black] at (1.35,-0.15) {$d_{11}$};
  \fill (0,0) circle (1.4pt);
\end{tikzpicture}
\caption{Вектор обратной решётки $\vb H$ перпендикулярен семейству узловых
плоскостей (здесь~— узловых прямых $(11)$ косоугольной двумерной решётки),
а его длина обратна расстоянию $d$ между ними}
\label{fig:H-vector}
\end{figure}

\paragraph{Закон зон (соотношение Вейсса)}
Условие $hu+kv+lw=0$ означает, что направление $[uvw]$ лежит в плоскости
$(hkl)$, т.\,е. плоскость принадлежит \emph{зоне} с осью $[uvw]$. Отсюда
практическое правило: ось зоны двух плоскостей $(h_1k_1l_1)$ и
$(h_2k_2l_2)$~— векторное произведение их индексов,
$[uvw]=[k_1l_2-l_1k_2,\ l_1h_2-h_1l_2,\ h_1k_2-k_1h_2]$, и аналогично
плоскость, содержащая два направления.

\begin{primer}[Ось зоны]
Плоскости $(110)$ и $(1\br10)$: $[uvw]=[1\cdot0-0\cdot(-1),\ 0\cdot1-1\cdot0,\
1\cdot(-1)-1\cdot1]=[00\br2]\parallel[001]$. Проверка: $1\cdot0+1\cdot0+0\cdot1=0$
и $1\cdot0-1\cdot0+0\cdot1=0$~— обе плоскости содержат $[001]$.
\end{primer}

\paragraph{Точки первого и второго рода}
В дифракционных методах (например, в методе Лауэ) полюсы плоскостей
называют точками первого рода, выходы направлений~— точками второго
рода. В кубических кристаллах они совпадают ($[hkl]\perp(hkl)$). В
тетрагональных совпадают лишь $\{hk0\}$ с $\nfam{hk0}$ и $\{001\}$ с
$\nfam{001}$; в гексагональных~— $\{hki0\}$ с $\nfam{hki0}$ и $\{0001\}$ с
$\nfam{0001}$ (в четырёхиндексной записи).

\paragraph{Узлы обратной решётки реального кристалла}
У идеального монокристалла узел обратной решётки~— точка. У реальных
объектов он размыт (рис.~\ref{fig:recip-blur}):
\begin{itemize}
\item в \emph{мозаичном} монокристалле (блоки слегка разориентированы)~—
по поверхности сферы радиуса $|\vb H|$, т.\,е. в направлениях,
перпендикулярных $\vb H$;
\item в идеальном \emph{поликристалле} со случайной ориентацией зёрен узел
превращается в сферу радиуса $|\vb H|$;
\item при \emph{неоднородной упругой деформации} межплоскостные расстояния
различны, и узел размывается вдоль $\vb H$.
\end{itemize}

\begin{figure}[ht]
\centering
\foreach \typ/\lab in {1/а, 2/б, 3/в}{%
\begin{tikzpicture}[scale=1.0]
  \fill (0,0) circle (1.4pt) node[below left=-1pt, font=\small]{$O$};
  \draw[fmGray, dashed] (0,0) circle (1.6);
  \ifnum\typ=1
    \fill[fmBlue!40, rotate=50] (1.6,0) ellipse (0.08 and 0.35);
  \fi
  \ifnum\typ=2
    \draw[fmBlue!60, line width=2.2pt] (0,0) circle (1.6);
  \fi
  \ifnum\typ=3
    \fill[fmBlue!40, rotate=50] (1.6,0) ellipse (0.35 and 0.08);
  \fi
  \draw[bvec, fmRed] (0,0) -- (50:1.55) node[midway, left, font=\small]{$\vb H$};
  \node[font=\small] at (0,-2.0) {\textit{\lab}};
\end{tikzpicture}\hspace{6mm}}
\caption{Размытие узла обратной решётки: \textit{а}~— мозаичный
монокристалл (вдоль сферы радиуса $|\vb H|$); \textit{б}~— поликристалл
(сфера); \textit{в}~— монокристалл с неоднородной упругой деформацией (вдоль
$\vb H$)}
\label{fig:recip-blur}
\end{figure}

\begin{podrobnee}[Обратное пространство как фурье-пространство]
Множество узлов можно задать решётчатой функцией~— суммой
$\delta$-функций в узлах:
$\text{Ш}(\vb r)=\sum_{m,n,p}\delta(\vb r-\vb R_{mnp})$ в прямом пространстве и
$\text{Ш}(\vb h)=\sum_{h,k,l}\delta(\vb h-\vb H_{hkl})$ в обратном. Эти функции
связаны преобразованием Фурье, поэтому обратное пространство называют
фурье-пространством. При дифракции рентгеновских лучей в прямом
пространстве задана электронная плотность $\rho(\vb r)$, а в обратном~—
амплитуда рассеяния $A(\vb h)$: обратное пространство~— это
\emph{дифракционное} пространство.
\end{podrobnee}
